% "Where the Money Comes From": shared by the rulebook and the platform manuals (docs/manual/*).
% Each document sets its own heading above this.
New \Vp{player} often feel busy without knowing what they are building toward. It helps to know where the
\Game{currency} actually comes from.

\begin{note}{IN ONE LINE}
\VCp{exploit} don't make money; \V{coin} does. You get paid when a \V{hive} you lead is \V{absorbed}, or when the game ends.
Running \Vp{exploit} is how you decide which \Vp{hive} get \V{absorbed}, and when.
\end{note}

\subsection{Founding is about position, not money}
An \V{infiltration} earns only the \V{finderscut}: one free \V{coin}, worth \Y{200}--\Y{400} while the \V{hive} is small. Its real
value is that you choose which syndicate comes up, and you are first to buy its \V{coin}, at the lowest price it will ever have.
Founding is the cheapest way to become \V{root} of a \V{hive}.

\subsection{The big money: \V{root} and \V{sudo}}
\begin{itemize}
  \item \textbf{The \V{absorbed} \V{hive} pays, not the \V{absorber}.} This is the part new \Vp{player} get backwards. At a \V{hivejack},
    only the syndicate being swallowed pays its \V{root} and \V{sudo} bonuses.
  \item \textbf{The bonuses are large.} \VC{root} pays ten times the \V{coin}'s price and \V{sudo} five times. At a price of
    \Y{600}, \V{root} is \Y{6,000}, as much as your whole starting stake.
  \item \textbf{So the classic play} is to lead a small or medium \V{hive} that sits next to a bigger one, then run the \V{exploit}
    that joins them, on your terms.
\end{itemize}

\subsection{The rest: the \V{coin} itself}
\begin{itemize}
  \item \textbf{At a \V{hivejack}}, the \V{absorbed} \V{coin} can be cashed out, swapped 2 for 1 into the \V{absorber}, or kept.
    A \V{swap} is how you climb into a big \V{hive} without spending \Game{currency}.
  \item \textbf{At the \V{nullroute}}, every \V{hive} still on the \V{grid} pays \V{root} and \V{sudo}, and all \V{coin} is sold at its
    final price. A \V{safe} \V{hive} can never be \V{absorbed}, so it pays only now, but by now it is usually big.
\end{itemize}

\subsection{Why it is hard to see: cash is the real limit}
You start with \Y{\Rule{startcash}} and may buy only \Rule{maxbuy} \Vp{coin} a turn, and nothing you hold pays until a \V{hivejack} or
the end. An early bonus counts twice: it pays you, and it pays for the \V{coin} that wins the next fight for \V{root}.

\begin{example}[: FINDER'S CUT VS.\ \MakeUppercase{\V{root}}]
Early on you found a low-tier \V{hive}: your \V{finderscut} is worth \Y{200}. You keep buying, and when the \V{hive} has grown
to 5 \Vp{node} you hold 4 of its \V{coin} and nobody else holds more than 2. A larger \V{hive} sits one \V{node} away, and you
hold that \V{exploit}. You run it: \Y{5,000} for \V{root}, then you \V{cashout} your 4 \V{coin} at \Y{500} each for \Y{2,000} more.
The \V{infiltration} was the doorway; the \V{hivejack} was the payday.
\end{example}
